\section{Sharp non-identification boundaries}\label{sec:nonid}

The positive results rely on structural invariance, not merely analytic regularity. We now show what fails when key modules are unrestricted.

Let a general mechanism have drift
\[
 AU_t+G(t,U_t).
\]
For bounded $L$, define
\[
 \mathcal T_L(A,G)=(A-L,G+L\,\cdot).
\]

\begin{theorem}[Operator--forcing gauge]\label{thm:gauge-main}
The mechanisms $(A,G)$ and $(A-L,G+L\,\cdot)$ generate the same factual stochastic evolution equation and hence the same factual path under common exogenous noise. Under a common operator intervention $A\mapsto\widetilde A$, their post-intervention drifts differ by $Lu$. Unless the selected target annihilates this direction, the post-intervention laws differ.
\end{theorem}

This is not merely an adapted-feedback pathology. Let $a_1\ne a_2$, let $h(0)=0$, and define an open-loop input
\[
 F_i(t)=Z(h'(t)+a_ih(t)),
\]
where $Z$ is generated ex ante. Both equations
\[
 dX_t=-a_iX_tdt+F_i(t)dt
\]
have the same complete factual path $X_t=Zh(t)$. After a common intervention to $\widetilde a$,
\[
 X_t^{i,do}
 =Z\left[
 h(t)+(a_i-\widetilde a)\int_0^te^{-\widetilde a(t-s)}h(s)ds
 \right],
\]
so the target laws separate.

A positive-noise construction works on every finite observation grid. With a piecewise-constant Gaussian input, the observed vector is
\[
 Y=L_aZ+\varepsilon_a.
\]
For a common covariance $C$, choosing
\[
 K_a=L_a^{-1}(C-C_a^W)L_a^{-\top}
\]
produces the same factual $N(0,C)$ law under distinct operators, while a common post-intervention operator generally produces different covariance matrices.

\begin{theorem}[Sharp boundary assembly]\label{thm:nonid-assembly}
Each of the following enlarged model classes contains observationally equivalent factual mechanisms with different post-intervention target laws:
\begin{enumerate}[label=(\roman*)]
\item unrestricted bounded state-feedback forcing;
\item unrestricted open-loop forcing law, even when independent of the Wiener noise;
\item unknown random initial law under finite observation;
\item rank-deficient sensing with a free target-relevant latent direction;
\item unknown calibration for field-scale targets;
\item intervention-dependent measurement;
\item non-identified state and measurement noise decomposition.
\end{enumerate}
Furthermore, identified marginal counterfactual laws do not identify a unit-level contrast without a cross-world coupling convention.
\end{theorem}

\begin{figure}[t]
\centering
\resizebox{.98\textwidth}{!}{\begin{tikzpicture}[font=\small]
\matrix (m) [matrix of nodes,
  nodes={minimum height=.62cm, anchor=center},
  column 1/.style={nodes={align=left, anchor=west, font=\small, text width=5.6cm}},
  column 2/.style={nodes={text width=1.7cm, align=center}},
  column 3/.style={nodes={text width=2.3cm, align=center}},
  column 4/.style={nodes={text width=2.1cm, align=center}},
  column 5/.style={nodes={text width=2.1cm, align=center}},
  row 1/.style={nodes={font=\footnotesize\bfseries, align=center}},
  row sep=1pt, column sep=4pt]
{
 |[font=\footnotesize\bfseries, align=left]| freed module (nuisance class) & factual law & fixed-interv.\ marginal $F_{a,w}$ & distributional effect $\Delta_{a,w,z}$ & unit-level contrast \\
 baseline: all modules invariant (Ass.~2.1, \S4) & \yes & \yes & \yes & \yes$^{\dagger}$ \\
 state-feedback forcing $G+L\,\cdot$ (gauge, Thm~8.1) & \yes & \no & \no & \no \\
 open-loop forcing law (ex-ante $Z$) & \yes & \no & \no & \no \\
 initial law free, finite observation & \yes & \no & \no & \no \\
 rank-deficient sensing, free latent direction & \yes & \no & \no & \no \\
 rank-deficient sensing, known cross-mode link & \yes & \yes & \yes & \yes$^{\dagger}$ \\
 unknown calibration $c$ (sensor $Y=c\,U$) & \yes & \half & \half & \no \\
 intervention-dependent measurement & \yes & \no & \no & \no \\
 cross-world coupling free & \yes & \yes & \yes & \no \\
};
% header rule and separators
\draw[black!60, thick] (m-1-1.south west) -- (m-1-5.south east);
\draw[black!25] (m-2-1.south west) -- (m-2-5.south east);
\draw[black!25] (m-1-1.north east) -- (m-10-1.south east);

% notes
\node[anchor=north west, align=left, font=\footnotesize, text=black!70, text width=14.2cm]
  at ($(m-10-1.south west)+(0,-0.25)$)
  {\half\ target-specific: field-scale targets fail, sensor-scale targets remain identified (Thm~8.2(v)).\quad
   $^{\dagger}$ unit-level contrasts additionally require the cross-world coupling convention (shared Wiener path, \S3);
   marginals never do.\\
   Reading: each row frees one module of the SCM while preserving observational equivalence of the factual law
   (column 1 always \yes); columns record which post-intervention targets survive (\S4, Thm~8.1--8.2).};
\end{tikzpicture}
}
\caption{Identification boundary map. Each row frees one structural module while preserving observational equivalence of the factual law; columns record which post-intervention targets survive (\Cref{thm:gauge-main,thm:nonid-assembly}).}
\label{fig:boundary}
\end{figure}

The qualifications are target specific. Unknown calibration can leave sensor-scale counterfactuals identified while field-scale targets fail. Partial sensing can be harmless when known cross-mode structure links latent and observed directions. Marginal CDF and quantile effects do not require identification of cross-world correlation.

The shared-Wiener construction of \Cref{sec:semantics} selects a structural coupling for unit-level contrasts. Other couplings can preserve the same marginals and change
\[
 \Var(X^a-X^b)=v_a+v_b-2\rho\sqrt{v_av_b}.
\]
This limits individual-effect claims but does not weaken the marginal distribution results.
